% !PRINTED 629
\weberpage{629}
%<<CONT
und wenn wir dies auf jeden Factor von $F_n(t)$ anwenden, und
beachten, dass $p^{h+1}$ nach dem Modul $m$ dieselbe Zahlenreihe
durchläuft, wie $p^h$, so folgt:
\[
[F_n(t)]^p \equiv F_n(t^p) \ (\mathrm{mod}\ p),
\]
und diese Congruenz gilt dann natürlich auch für den Modul $\mathfrak p$.

Dies ist aber das Kennzeichen dafür, dass $F_n(t)$ mit einem
Polynom $P_n(t)$ mit ganzen rationalen Zahlenco\"efficienten nach
dem Modul $\mathfrak p$ congruent ist (\S.\,150, 4.), also
\weq{4}{F_n(t) \equiv P_n(t) \ (\mathrm{mod}\ \mathfrak p).}

Setzen wir dies in (3) ein, so ergiebt sich zunächst eine
Congruenz nach dem Modul $\mathfrak p$, die aber, da es eine Congruenz
zwischen \gesperrt{rationalen} Functionen ist, auch für den Modul $p$
bestehen muss, also
\weq{5}{f(t) \equiv P_{\xi_1}(t)\ P_{\xi_2}(t)\ \ldots\ P_{\xi_e}(t) \ (\mathrm{mod}\ p).}

Nach der Definition (1), (4) ist
\weq{6}{P_1(r) \equiv 0 \ (\mathrm{mod}\ \mathfrak p),}
und die sämmtlichen Wurzeln dieser Congruenz sind
\weq{7}{r,\ r^p,\ \ldots,\ r^{p^{f-1}}.}

Ebenso hat die Congruenz
\[
P_n(t) \equiv 0 \ (\mathrm{mod}\ \mathfrak p)
\]
die Wurzeln
\[
r^n,\ r^{np},\ \ldots,\ r^{np^{f-1}},
\]
und diese Wurzeln sind, wenn man $n$ das System der Zahlen
$\xi_1, \xi_2, \ldots, \xi_e$ durchlaufen lässt, alle incongruent nach dem
Modul $\mathfrak p$.

Macht man in der Congruenz
\[
F_1(t) \equiv P_1(t) \ (\mathrm{mod}\ \mathfrak p)
\]
die Substitution $(r, r^n)$, so geht $F_1(t)$ in $F_n(t)$, $\mathfrak p$ in das
conjugirte Ideal $\mathfrak p_n$ über, und $P_1(t)$ bleibt als rationale Form
ungeändert. Wir haben also
\weq{8}{F_n(t) \equiv P_1(t) \ (\mathrm{mod}\ \mathfrak p_n),}
und wenn wir hierin $t=r$ setzen, so folgt:
\weq{9}{F_n(r) \equiv P_1(r) \ (\mathrm{mod}\ \mathfrak p_n),}
$F_n(r)$ ist aber, wenn $n$ nicht in $\mathfrak A$ enthalten ist, relativ prim
zu $p$, also nicht durch $\mathfrak p_n$ theilbar.

Es ist also $P_1(r)$ nach (9) durch $\mathfrak p_1$, aber durch keinen der
%CONT>>
